% Content of the QECLean blueprint, shared by the web and pdf versions.
%
% The nodes themselves are NOT written here. They are extracted from Lean by
% LeanArchitect out of the annotations in QECBlueprint.lean, and pulled in by
% the \input below; this file supplies the narrative and decides the order in
% which nodes appear. To change a statement or a proof sketch, edit the
% annotation in QECBlueprint.lean and re-run `lake build QECBlueprint:blueprint`
% -- do not edit the generated files under .lake/build/blueprint.

\input{../../.lake/build/blueprint/library/QECBlueprint}

\chapter{Introduction}
\label{chap:intro}

This is the blueprint for \href{https://github.com/Stavan-Jain/QECLean}{QECLean},
a formalization of the stabilizer formalism for quantum error correction in Lean~4
and Mathlib.

A blueprint is a map of a formalization: informal mathematical exposition, one
node per definition or theorem, wired into a dependency graph that shows how the
results rest on one another and how much of the plan is already formal. Here
every node is \emph{already} formal --- the library is complete and sorry-free ---
so the graph is best read as a guide to the architecture rather than as a
to-do list.

\section{What is proved}

Two results anchor the development.

\begin{itemize}
  \item \textbf{The toric family.} For every $L \ge 2$, the $L \times L$ toric
    code is verified as an $[[2L^2, 2, L]]$ stabilizer code: the distance is
    exactly $L$, not merely bounded below by it
    (\cref{thm:toric-distance}).
  \item \textbf{The gross code.} IBM's $[[144,12,12]]$ bivariate-bicycle code
    --- the flagship qLDPC code of Bravyi et al. --- is verified with distance
    exactly $12$, unconditionally (\cref{thm:gross-distance}).
\end{itemize}

Alongside these sit the rotated surface family, the repetition codes, the
iceberg family, and the small named codes: Steane's $[[7,1,3]]$, Shor's
$[[9,1,3]]$, the perfect $[[5,1,3]]$ code, and the $[[4,2,2]]$ detection code
(\cref{chap:codes}).

\section{The axiom discipline}

Everything on the main branch depends on exactly Mathlib's three standard
axioms --- \texttt{propext}, \texttt{Classical.choice}, \texttt{Quot.sound} ---
and nothing else. In particular there is no \texttt{native\_decide}, which
would replace a kernel check by a compiler-trust axiom, and no \texttt{sorry}.
An input that cannot yet be proved is carried as a named hypothesis on the
theorem, where it is visible in the statement rather than hidden in the axiom
set. This is enforced in CI, so a violation fails the build rather than landing
quietly.

The discipline is not free. The gross code distance proof was originally
carried by 135 \texttt{native\_decide} leaves and had to be converted to
kernel-checkable form: packed natural-number tables instead of array lookups,
quantifier bridges so the kernel enumerates concrete lambdas, sparse rewrites
in place of large finite sums, and Gaussian-style certificates where a sweep
would otherwise be needed.

\section{Reading the graph}

An edge $A \to B$ in the dependency graph means that $B$ is the nearest
blueprint ancestor of $A$: LeanArchitect walks the type and the proof term of
each annotated declaration and stops as soon as it reaches another annotated
declaration. The edges are therefore read off the real proof terms rather than
maintained by hand, and a node that a proof genuinely stopped needing loses its
edge as soon as the proof changes.

\chapter{The Pauli group}
\label{chap:pauli}

Everything begins with the group $\mathcal{P}_n$ of $n$-qubit Pauli operators
with phases. The formalization splits an element into a phase exponent in
$\Z_4$ and a phase-free operator $\mathrm{Fin}\,n \to \{I,X,Y,Z\}$. The split
is what makes multiplication computable: the operator part multiplies
qubitwise, and the phases that fall out of products such as $XZ = -iY$
accumulate separately.

\inputleannode{def:pauli-operator}
\inputleannode{def:pauli-group-element}
\inputleannode{def:nqubit-pauli-operator}
\inputleannode{def:nqubit-pauli-group}

The syntactic representation is tied to actual operators on a Hilbert space by
a matrix interpretation; without it the development would be a study of a
finite group with no quantum content.

\inputleannode{def:pauli-to-matrix}
\inputleannode{def:pauli-mulop}
\inputleannode{thm:pauli-mul-assoc}

\section{Weight}

Code distance is a statement about how many qubits an operator touches, so
support and weight are the measures everything else is phrased in.

\inputleannode{def:pauli-support}
\inputleannode{def:pauli-weight}

\section{Commutation}

The single most used fact in the development is that commutation in
$\mathcal{P}_n$ is a parity condition. Two Paulis pick up a sign $-1$ for each
qubit where they disagree nontrivially, so the global sign is $(-1)^k$ and
there is no third possibility between commuting and anticommuting.

\inputleannode{def:anticommutes-at}
\inputleannode{thm:commutes-iff-even}
\inputleannode{def:anticommute}
\inputleannode{thm:commute-or-anticommute}

\chapter{The binary symplectic representation}
\label{chap:symplectic}

Parity is linear algebra over $\F_2$, and that observation is the engine of the
whole subject. Forgetting phases sends $\mathcal{P}_n$ onto $\F_2^{2n}$,
products become sums, and --- by \cref{thm:commutes-iff-symplectic} ---
commutation becomes orthogonality under a symplectic form. An abelian subgroup
becomes a self-orthogonal subspace, and independence of generators becomes a
rank computation.

\inputleannode{def:to-symplectic}
\inputleannode{thm:to-symplectic-injective}
\inputleannode{def:symplectic-inner}
\inputleannode{thm:commutes-iff-symplectic}
\inputleannode{def:check-matrix}
\inputleannode{def:rows-linear-independent}

\chapter{Stabilizer groups and the codespace}
\label{chap:stabilizer}

\inputleannode{def:is-stabilized-by}
\inputleannode{def:stabilizer-group}
\inputleannode{thm:neg-identity-not-mem}
\inputleannode{def:codespace}

The two conditions in \cref{def:stabilizer-group} --- abelian, and not
containing $-I$ --- are exactly what is needed for the codespace to be nonzero.
Commutativity makes the eigenspace projectors compatible; excluding $-I$ rules
out the contradiction $\psi = -\psi$. The proof that the codespace is genuinely
nonempty runs through the group sum.

\inputleannode{def:stabilizer-sum}
\inputleannode{thm:codespace-nonempty}

\section{The centralizer}

Operators that preserve the codespace are those commuting with every
stabilizer. For the Pauli group --- and this is special to it --- normalizing
and centralizing coincide, which is why the literature uses the two words
interchangeably in this setting.

\inputleannode{def:centralizer}
\inputleannode{thm:normalizer-eq-centralizer}
\inputleannode{thm:stabilizer-le-centralizer}

\chapter{Logical operators, codes, and distance}
\label{chap:codes-abstract}

The encoded information is acted on by the quotient
$\mathcal{C}(\mathcal{S})/\mathcal{S}$: operators that preserve the codespace,
modulo those that act trivially on it.

\inputleannode{def:pauli-logical}
\inputleannode{thm:logical-iff-generators}
\inputleannode{def:nontrivial-logical}
\inputleannode{thm:nontrivial-logical-iff}

\Cref{def:nontrivial-logical} has three conditions, not two, and the third is
easy to forget: no stabilizer may share the operator part of $g$. Without it a
$g$ differing from a stabilizer only by a phase would count as acting
nontrivially, and the CSS bridge argument of \cref{thm:not-both-boundary}
would fail.

\inputleannode{thm:nontrivial-logical-transfer}
\inputleannode{def:logical-qubit-ops}

\section{Packaging a code}

\inputleannode{def:generators-independent}
\inputleannode{thm:independent-of-check-matrix}
\inputleannode{def:stabilizer-code}

Instantiating \cref{def:stabilizer-code} for a concrete family \emph{is} the
theorem that the family encodes $k$ qubits into $n$: the type carries $n$ and
$k$, and the obligations --- generator count, independence, pairwise
commutation, exclusion of $-I$, and the centralizer and anticommutation
conditions on the logical operators --- must all be discharged to build an
inhabitant.

\section{Distance}

\inputleannode{def:has-code-distance}
\inputleannode{thm:distance-min-weight}
\inputleannode{def:code-with-distance}

Both halves of \cref{def:has-code-distance} matter. The lower bound is the
error-correction guarantee; the witness of exact weight is what makes the
distance a value rather than a bound, and it is usually the easier half to
forget.

\chapter{CSS structure}
\label{chap:css}

A CSS code is one whose stabilizer generators are each purely $X$-type or
purely $Z$-type. The separation is what lets a homological argument apply: the
$X$ and $Z$ sides become the two ends of a chain complex.

\inputleannode{def:x-type}
\inputleannode{def:z-type}
\inputleannode{thm:css-distance-two}
\inputleannode{thm:css-distance-three}

\chapter{The homological framework}
\label{chap:homological}

This chapter is the abstract heart of the development. A length-three chain
complex over $\F_2$ induces a CSS code: $1$-chains are the physical qubits,
$\partial_1 \partial_2 = 0$ is the commutation of the $X$- and $Z$-checks, and
--- the slogan --- \emph{logical operators are homology classes}. Both the
toric family and the gross code are instances, and both inherit their distance
argument from \cref{thm:chain-weight-transfer}.

\inputleannode{def:homological-code}
\inputleannode{thm:boundary-comp-zero}
\inputleannode{def:cycles}
\inputleannode{def:boundaries}
\inputleannode{thm:boundaries-le-cycles}
\inputleannode{def:H1}
\inputleannode{thm:finrank-H1}

\section{From chains to Paulis}

\inputleannode{def:chain-x-operator}
\inputleannode{def:chain-z-operator}
\inputleannode{thm:chain-x-add}
\inputleannode{thm:x-logical-iff-cycle}

\section{Distance from homology}

\inputleannode{thm:not-both-boundary}
\inputleannode{thm:chain-weight-transfer}

\Cref{thm:chain-weight-transfer} is the load-bearing step: it is what allows a
statement about minimum weights of homology representatives --- a question in
$\F_2$ linear algebra --- to settle a statement about the weights of Pauli
operators. Every concrete distance theorem below is an application of it.

\chapter{The toric code}
\label{chap:toric}

The $L \times L$ square lattice on the torus, read as a chain complex. The
$2L^2$ edges are the qubits; vertices carry the $Z$-checks and faces the
$X$-checks.

\inputleannode{def:toric-chain-complex}
\inputleannode{def:toric-H1}
\inputleannode{thm:toric-H1-dim}

\section{Wrapping invariants}

The homology of the torus is detected by two $\F_2$-valued winding numbers,
counting how often a chain wraps horizontally and vertically. They are defined
on all chains, are unchanged by adding a boundary, and together separate
homology classes.

\inputleannode{def:wrapping-h}
\inputleannode{def:wrapping-v}
\inputleannode{thm:wrapping-boundary-zero}
\inputleannode{def:toric-phi}
\inputleannode{thm:toric-phi-equiv}

\section{The code and its distance}

\inputleannode{def:toric-code}

The natural generating set of the toric code has $2L^2$ elements, two more than
the $n - k = 2L^2 - 2$ that \cref{def:stabilizer-code} demands: the vertex
checks multiply to the identity, and so do the face checks. The packaging
therefore uses a trimmed generator list, and the homological identities are
what prove the trimmed list generates the same subgroup. Distance proofs stated
against the untrimmed stabilizer group are carried across by
\cref{thm:nontrivial-logical-transfer}.

\inputleannode{thm:toric-dX}
\inputleannode{thm:toric-dZ}
\inputleannode{thm:toric-distance}
\inputleannode{def:toric-code-with-distance}

\chapter{Bivariate-bicycle codes and the gross code}
\label{chap:bb}

Bivariate-bicycle codes are qLDPC codes built from two commuting polynomial
shift operators on an abelian group. The gross code is the $[[144,12,12]]$
member of the family; the argument here derives its distance from that of its
$[[72,12,6]]$ base by exhibiting the gross complex as a free $\Z_2$ cover of
the base and showing the distance doubles.

\inputleannode{def:x-double-cover}

\section{The base floor}

The base contributes a floor of six, established by ruling out every light
cycle with a finite certificate that the kernel checks directly.

\inputleannode{def:small-cycle-data}
\inputleannode{thm:small-cycle-floor}
\inputleannode{thm:base-distance}
\inputleannode{thm:gross-logical-weight}

\section{Doubling}

Pushing a gross logical down to the base gives weight at least six. Getting to
twelve requires knowing when the two sheets of the cover contribute
independently. They do outside a \emph{dangerous} sector, where the deck
transformation acts trivially on homology; there the light stabilizers are
classified explicitly and the surviving classes are confined to Smith cosets
whose floor is again twelve.

\inputleannode{thm:gross-sector-dichotomy}
\inputleannode{def:gross-stabilizer-code}
\inputleannode{thm:gross-distance}
\inputleannode{def:gross-code-with-distance}

\chapter{Code instances}
\label{chap:codes}

The remaining codes exercise the same machinery on smaller or differently
shaped examples. The $[[5,1,3]]$ code is the interesting outlier: its
generators mix $X$ and $Z$ on the same qubit, so it is not CSS, the homological
route is unavailable, and its distance is settled instead by a direct search
over weight-one and weight-two anti-witnesses.

\inputleannode{def:steane7}
\inputleannode{def:shor9}
\inputleannode{def:five-qubit}
\inputleannode{def:four-qubit}
\inputleannode{def:iceberg}
\inputleannode{def:rotated-surface}
\inputleannode{def:repetition-n}
